% 復習5　問いを磨く
\session{5}{1}{60分}{問いを磨く}{AIによる批判的検証}{}

\goals{
  \item AIを「厳しい査読者」として使い、問いの弱点を洗い出せる
  \item 指摘に対して、補強するか問いを修正するかを、根拠をもって決められる
  \item 問い・背景・先行研究・調べ方の見通しを1枚にまとめられる
}

\work{準備}{AIに提示する材料}
\answerbox[仮の問い]{14mm}
\answerbox[背景（200字程度）]{28mm}
\answerbox[主な先行研究（3件）]{20mm}

\work[70mm]{ワーク1・2}{査読者の指摘に、補強か修正で答える}
\begin{promptbox}[査読者として批判させる]
あなたは厳しい査読者です。次のリサーチ・クエスチョンについて、1. 数か月で調べて答えられるか　2. 先行研究に対する新しさがあるか　3. 範囲は適切か　4. 用語は明確か　の観点から問題点を指摘し、それぞれ改善案を示してください。問い：（あなたの問い）　背景：（200字程度）　主な先行研究：（3件）
\end{promptbox}
\begin{wstabh}{colspec={Q[l,wd=22mm,m] X[3,l,m] Q[c,wd=14mm,m] X[3,l,m]},row{2-Z}={ht=17mm}}
  観点 & AIの指摘（要点） & 補強／\newline 修正 & 自分の対応と、その根拠 \\
  答えられるか & & & \\ 新しさ & & & \\ 範囲 & & & \\ 用語の明確さ & & & \\
\end{wstabh}

\begin{promptbox}[反論に答える練習]
私の問いに対して「（AIが挙げた反論）」という批判があります。この批判に答える論点を3つ挙げ、それぞれどんな根拠が必要かを示してください。根拠が見つかっていない論点は、そのように明示してください。
\end{promptbox}

\work{修正の記録}{問いはどう変わったか}
\twoboxes{修正前の問い}{修正後の問い（確定）}{24mm}

\work[90mm]{ワーク3}{確定した問いを1枚にまとめる}
\begin{wstab}{colspec={Q[l,wd=30mm,m] X[l,m]},row{1}={ht=14mm},row{2-Z}={ht=22mm}}
  \hc{問い} & \\
  \hc{背景と意義} & \\
  \hc{先行研究\newline でわかって\newline いること} & \\
  \hc{調べ方の見通し} & \\
\end{wstab}

\work[40mm]{振り返り}{復習1〜5のチェック}
\begin{checks}
  \item AIが挙げた文献・数値を、そのまま使わずに確かめた
  \item 問いの候補から、自分で基準を決めて選んだ
  \item 文献リストの必須項目をすべて揃えた
  \item AIの要約を原文と照らし合わせた
  \item AIの指摘を受けて、問いをどう変えたかを記録した
  \item AI活用記録を毎回書いた
\end{checks}
\aimemo
